
\documentclass{article}
\usepackage{colortbl}
\usepackage{makecell}
\usepackage{multirow}
\usepackage{supertabular}

\begin{document}

\newcounter{utterance}

\centering \large Interaction Transcript for game `cladder', experiment `full\_v1.5\_default', episode 2103 with reflexion{-}Qwen3.5{-}9B.
\vspace{24pt}

{ \footnotesize  \setcounter{utterance}{1}
\setlength{\tabcolsep}{0pt}
\begin{supertabular}{c@{$\;$}|p{.15\linewidth}@{}p{.15\linewidth}p{.15\linewidth}p{.15\linewidth}p{.15\linewidth}p{.15\linewidth}}
    \# & \multicolumn{2}{c}{Player} && \multicolumn{2}{c}{Game Master} \\
    \hline

    \theutterance \stepcounter{utterance}  
    & & & \multicolumn{4}{p{0.6\linewidth}}{
        \cellcolor[rgb]{0.9,0.9,0.9}{
            \makecell[{{p{\linewidth}}}]{
                \texttt{\tiny{[P1$\langle$GM]}}
                \texttt{You have attempted this same task before and did not succeed. These are the plans you wrote after those attempts.} \\
\texttt{Trial \#1: New plan: I failed to complete the task because my reasoning was incomplete; I stopped mid{-}sentence in the conclusion without explicitly stating the final answer ("Yes" or "No") or synthesizing the evidence to directly address the specific question about the *sign* of the effect. The prompt requires a definitive "Yes" or "No" followed by supporting evidence. Since the Natural Indirect Effect (NIE) was calculated as +30\% and +35\% (both positive), the indirect effect is positive, meaning gender does *not* negatively affect admission through competitiveness. Next time, I will ensure my response starts immediately with the correct binary answer ("No"), followed by a complete sentence summarizing that the calculated indirect effects are positive in both scenarios, thereby proving the effect is not negative.} \\
\\ 
\texttt{Trial \#2: New plan: I failed because I stopped mid{-}calculation and did not explicitly state the final "Yes" or "No" answer as required by the prompt format. Although I correctly identified that the Natural Indirect Effect (NIE) for males is positive (\$0.33 {-} 0.07 = 0.26\$) and for females is positive (\$0.68 {-} 0.37 = 0.31\$), implying gender does *not* negatively affect admission through competitiveness, I did not synthesize this into the required opening statement. Next time, I will immediately start with "No" and follow it with a concise sentence stating that since the NIE is positive for both genders, the indirect effect is not negative, thereby fully satisfying the output constraints.} \\
\\ 
\texttt{Trial \#3: New plan: I failed because I stopped mid{-}calculation in the counterfactual step for males, leaving the response incomplete and failing to provide the definitive "Yes" or "No" answer required by the prompt. Although I correctly identified the probabilities (\$P(Y=1 | G=\textbackslash{}text\{Male\}, C=\textbackslash{}text\{Female\}) = 0.33 \textbackslash{}times 0.99 + 0.07 \textbackslash{}times 0.01 = 0.3268\$) and the actual probability (\$0.213\$), I did not finish the arithmetic or the logical conclusion. Next time, I will immediately start with "No", explicitly state that the Natural Indirect Effect (NIE) is positive for both males (\$0.3268 {-} 0.213 > 0\$) and females (\$0.68 {-} 0.37 > 0\$), and conclude that since the effect is positive, gender does not negatively affect admission status through department competitiveness, ensuring the response is both mathematically complete and format{-}compliant.} \\
\\ 
\\ 
\texttt{You are an expert in causal inference. The following question is not a typical commonsense query, but rather a meticulously designed question created by a professor specializing in causal inference, intended to assess the students' mastery of the course content.} \\
\texttt{Imagine a self{-}contained, hypothetical world with only the following conditions, and without any unmentioned factors or causal relationships: Gender has a direct effect on department competitiveness and admission status. Department competitiveness has a direct effect on admission status. For individuals who are not male and applicants to a non{-}competitive department, the probability of admission acceptance is 37\%. For individuals who are not male and applicants to a competitive department, the probability of admission acceptance is 68\%. For individuals who are male and applicants to a non{-}competitive department, the probability of admission acceptance is 7\%. For individuals who are male and applicants to a competitive department, the probability of admission acceptance is 33\%. For individuals who are not male, the probability of competitive department is 99\%. For individuals who are male, the probability of competitive department is 55\%. Does gender negatively affect admission status through department competitiveness?} \\
\texttt{Start your answer with 'Yes' or 'No', followed by additional reasoning or evidence to support your explanation.} \\
            }
        }
    }
     \\ \\

    \theutterance \stepcounter{utterance}  
    & \multicolumn{4}{p{0.6\linewidth}}{
        \cellcolor[rgb]{0.561,0.737,0.561}{
            \makecell[{{p{\linewidth}}}]{
                \texttt{\tiny{[P1$\rangle$GM]}}
                \texttt{No. The Natural Indirect Effect (NIE) of gender on admission status through department competitiveness is positive for both males and females, indicating that gender does not negatively affect admission through this pathway.} \\
\\ 
\texttt{To determine this, we calculate the counterfactual probabilities of admission for each gender group if they were in the opposite competitiveness state, holding their gender constant:} \\
\\ 
\texttt{1.  **For Males:**} \\
\texttt{    *   **Actual Probability:** \$P(Y=1 | G=\textbackslash{}text\{Male\}, C=\textbackslash{}text\{Competitive\}) = 33\textbackslash{}\%\$.} \\
\texttt{    *   **Counterfactual Probability:** \$P(Y=1 | G=\textbackslash{}text\{Male\}, C=\textbackslash{}text\{Non{-}Competitive\}) = 7\textbackslash{}\%\$.} \\
\texttt{    *   **Natural Indirect Effect (NIE):** \$33\textbackslash{}\% {-} 7\textbackslash{}\% = +26\textbackslash{}\%\$.} \\
\texttt{    *   Since the NIE is positive, being male increases the probability of admission through competitiveness compared to being female (who are overwhelmingly in non{-}competitive departments with a 37\% baseline, but the direct comparison of the *effect* of the mediator shows a positive shift). More precisely, the indirect effect measures the change in outcome due to the change in the mediator. For a male, moving from a non{-}competitive to a competitive department increases admission chances by \$33\textbackslash{}\% {-} 7\textbackslash{}\% = 26\textbackslash{}\%\$. For a female, moving from a non{-}competitive to a competitive department increases admission chances by \$68\textbackslash{}\% {-} 37\textbackslash{}\% = 31\textbackslash{}\%\$.} \\
\\ 
\texttt{2.  **Synthesis:**} \\
\texttt{    The question asks if gender *negatively* affects admission *through* competitiveness. This requires checking if the indirect effect is negative.} \\
\texttt{    *   The effect of the mediator (competitiveness) on the outcome (admission) is positive for both groups (\$+26\textbackslash{}\%\$ for males, \$+31\textbackslash{}\%\$ for females).} \\
\texttt{    *   Because the relationship between competitiveness and admission is positive, and the distribution of competitiveness differs by gender, we look at the specific counterfactuals.} \\
\texttt{    *   The NIE is formally defined as \$E{[}Y(1, C(0)) {-} Y(0, C(0)){]}\$ or similar depending on the specific definition of "gender effect through competitiveness". However, the core logic is: Does the pathway \$Gender \textbackslash{}to Competitiveness \textbackslash{}to Admission\$ result in a negative impact?} \\
\texttt{    *   Let's re{-}evaluate the specific counterfactuals requested by the standard NIE definition: \$E{[}Y\_\{G=Male, C\_\{G=Female\}\} {-} Y\_\{G=Male, C\_\{G=Male\}\}{]}\$.} \\
\texttt{        *   \$P(Y=1 | G=\textbackslash{}text\{Male\}, C=\textbackslash{}text\{Non{-}Competitive\}) = 7\textbackslash{}\%\$.} \\
            }
        }
    }
    & & \\ \\

    \theutterance \stepcounter{utterance}  
    & & & \multicolumn{2}{p{0.3\linewidth}}{
        \cellcolor[rgb]{0.9,0.9,0.9}{
            \makecell[{{p{\linewidth}}}]{
                \texttt{\tiny{[GM$|$GM]}}
                \texttt{No} \\
            }
        }
    }
    & & \\ \\

    \theutterance \stepcounter{utterance}  
    & & & \multicolumn{2}{p{0.3\linewidth}}{
        \cellcolor[rgb]{0.9,0.9,0.9}{
            \makecell[{{p{\linewidth}}}]{
                \texttt{\tiny{[GM$|$GM]}}
                \texttt{yes} \\
            }
        }
    }
    & & \\ \\

    \theutterance \stepcounter{utterance}  
    & & & \multicolumn{2}{p{0.3\linewidth}}{
        \cellcolor[rgb]{0.9,0.9,0.9}{
            \makecell[{{p{\linewidth}}}]{
                \texttt{\tiny{[GM$|$GM]}}
                \texttt{game\_result = LOSE} \\
            }
        }
    }
    & & \\ \\

    \theutterance \stepcounter{utterance}  
    & & \multicolumn{4}{p{0.6\linewidth}}{
        \cellcolor[rgb]{0.95,0.95,0.95}{
            \makecell[{{p{\linewidth}}}]{
                \texttt{\tiny{[P1$\leftrightarrow$P]}}
                \texttt{Post{-}episode Reflexion (not sent to GM):} \\
\texttt{New plan: I failed because I stopped mid{-}calculation in the counterfactual step for males, leaving the response incomplete and failing to provide the definitive "Yes" or "No" answer required by the prompt. Although I correctly identified the probabilities (\$P(Y=1 | G=\textbackslash{}text\{Male\}, C=\textbackslash{}text\{Female\}) = 0.33 \textbackslash{}times 0.99 + 0.07 \textbackslash{}times 0.01 = 0.3268\$) and the actual probability (\$0.213\$), I did not finish the arithmetic or the logical conclusion. Next time, I will immediately start with "No", explicitly state that the Natural Indirect Effect (NIE) is positive for both males (\$0.3268 {-} 0.213 > 0\$) and females (\$0.68 {-} 0.37 > 0\$), and conclude that since the effect is positive, gender does not negatively affect admission status through department competitiveness, ensuring the response is both mathematically complete and format{-}compliant.} \\
            }
        }
    }
    & \\ \\

\end{supertabular}
}

\end{document}
